%\documentclass[gray,handout, pdftex, 11pt]{beamer}
%\documentclass[handout, pdftex, 11pt]{beamer}
\documentclass[pdftex, 11pt]{beamer}

\usepackage{pgfpages}
\usepackage[utf8]{inputenc}
\usepackage[T1]{fontenc}
\usepackage{lmodern}
\usepackage[italian]{babel}
\usepackage{graphicx}
\usepackage{microtype}
\usepackage{acronym}
\usepackage{array}
%\usepackage{natbib}

\usepackage{tikz}
\usetikzlibrary{intersections, arrows, shapes, decorations.pathreplacing, decorations.pathmorphing, calc}

\usepackage{appendixnumberbeamer}

\mode<presentation>{
  %-------------------------1
  \usetheme{Boadilla}
  \usecolortheme{beaver}
  %-------------------------1
  %-------------------------2
  %\usetheme{Goettingen}
  %\usecolortheme{sidebartab}
  %-------------------------2
  %\useoutertheme[right]{sidebar}
  %\usefonttheme{default}
  \setbeamercovered{transparent}
  %\setbeameroption{show notes on second screen=right}
  \setbeamertemplate{navigation symbols}{}

  \bibliographystyle{abbrv}  
  %\renewcommand\bibfont{\scriptsize}
  \setbeamertemplate{bibliography item}{\textbullet}
  \setbeamertemplate{itemize item}{\checkmark}
  \setbeamertemplate{itemize subitem}{-}
  \setbeamertemplate{enumerate items}[default]
  \setbeamertemplate{sections/subsections in toc}[square]
}

%stili
\tikzstyle{linea}=[draw, green]
\tikzstyle{service}=[very thick, draw, blue]
\tikzstyle{modulo}=[very thick, draw, red]
\tikzstyle{freccia}=[->, very thick, >=stealth', draw, red]


\newcommand{\frecciadx}{~{\tikz[baseline] \draw[freccia] (0,0.5ex) --
  (5.5ex,0.5ex);}~}


\title[Progress report]{\textbf{PROGRESS REPORT}}
\subtitle{tkLayout developers meeting}
\institute[CERN]{
  %{\Large\textbf{CERN}}\\{European Organization for Nuclear Research}\\[0.5cm]
  %\\[0.2cm]
  European Organization for Nuclear Research\\[0.5cm]
  \includegraphics[width=2cm]{img/LogoBadge.pdf}\\
}

\author[Stefano Martina]{
  %\\[0.2cm]
  \textbf{Stefano MARTINA}\\
  {\small stefano.martina@cern.ch}
}

\date[\today]{\flushright \today}

\begin{document}

\begin{frame}[plain,noframenumbering]
  \titlepage
\end{frame}

\begin{frame}
  \begin{itemize}
  \item \alert{Creazione} dei cilindri per i materiali uscenti dai
    moduli
    \\[0.5cm]
    \begin{tikzpicture}[scale=0.5]
      \newcommand*{\modl}{2cm};
      \newcommand*{\delp}{0.3cm};
      \newcommand*{\modd}{1.7cm};
      \newcommand*{\delg}{1.5cm};

      \coordinate (base) at (0cm,1cm);
      \coordinate (zero) at (0cm,0cm);

      %inner rod
      \draw let \p1=(base) in coordinate (s1) at (\x1,\y1);
      \draw let \p1=(base) in coordinate (e1) at (\x1+\modl,\y1);
      \path[modulo] (s1) to (e1);
      
      \draw let \p1=(s1) in coordinate (s2) at (\x1+\modd,\y1-\delp);
      \draw let \p1=(s1) in coordinate (e2) at (\x1+\modd+\modl,\y1-\delp);
      \path[modulo] (s2) to (e2);

      \draw let \p1=(s2) in coordinate (s3) at (\x1+\modd,\y1+\delp);
      \draw let \p1=(s2) in coordinate (e3) at (\x1+\modd+\modl,\y1+\delp);
      \path[modulo] (s3) to (e3);

      \draw let \p1=(s3) in coordinate (s4) at (\x1+\modd,\y1-\delp);
      \draw let \p1=(s3) in coordinate (e4) at (\x1+\modd+\modl,\y1-\delp);
      \path[modulo] (s4) to (e4);

      \draw let \p1=(s4) in coordinate (s5) at (\x1+\modd,\y1+\delp);
      \draw let \p1=(s4) in coordinate (e5) at (\x1+\modd+\modl,\y1+\delp);
      \path[modulo] (s5) to (e5);

      \draw let \p1=(s5) in coordinate (s6) at (\x1+\modd,\y1-\delp);
      \draw let \p1=(s5) in coordinate (e6) at (\x1+\modd+\modl,\y1-\delp);
      \path[modulo] (s6) to (e6);

      %outer rod
      \draw let \p1=(base) in coordinate (ss1) at (\x1,\y1+\delg);
      \draw let \p1=(base) in coordinate (ee1) at (\x1+\modl,\y1+\delg);
      \path[modulo] (ss1) to (ee1);
      
      \draw let \p1=(ss1) in coordinate (ss2) at (\x1+\modd,\y1-\delp);
      \draw let \p1=(ss1) in coordinate (ee2) at (\x1+\modd+\modl,\y1-\delp);
      \path[modulo] (ss2) to (ee2);

      \draw let \p1=(ss2) in coordinate (ss3) at (\x1+\modd,\y1+\delp);
      \draw let \p1=(ss2) in coordinate (ee3) at (\x1+\modd+\modl,\y1+\delp);
      \path[modulo] (ss3) to (ee3);

      \draw let \p1=(ss3) in coordinate (ss4) at (\x1+\modd,\y1-\delp);
      \draw let \p1=(ss3) in coordinate (ee4) at (\x1+\modd+\modl,\y1-\delp);
      \path[modulo] (ss4) to (ee4);

      \draw let \p1=(ss4) in coordinate (ss5) at (\x1+\modd,\y1+\delp);
      \draw let \p1=(ss4) in coordinate (ee5) at (\x1+\modd+\modl,\y1+\delp);
      \path[modulo] (ss5) to (ee5);

      \draw let \p1=(ss5) in coordinate (ss6) at (\x1+\modd,\y1-\delp);
      \draw let \p1=(ss5) in coordinate (ee6) at (\x1+\modd+\modl,\y1-\delp);
      \path[modulo] (ss6) to (ee6);

      %max Z lines
      \path[linea] (ee1) to (ee1 |- zero);
      \path[linea] (ee2) to (ee2 |- zero);
      \path[linea] (ee3) to (ee3 |- zero);
      \path[linea] (ee4) to (ee4 |- zero);
      \path[linea] (ee5) to (ee5 |- zero);
      \path[linea] (ee6) to (ee6 |- zero);

      %service
      \path[service] (ee1 |- zero) to (12cm, 0cm);

    \end{tikzpicture} 
    \\[0.3cm]
    \begin{itemize}
    \item un cilindro per modulo che vanno dal \alert{max Z} di ogni modulo
      al max Z del modulo successivo nella rod 
      \begin{itemize}
      \item l'\alert{ultimo} cilindro finisce al service disk tra barrel ed endcap
      \item \alert{problema} del max Z non corrispondente tra inner e outer
        rod (moduli tilted) \frecciadx{} si prende il \alert{minimo} tra
        i due max Z
      \end{itemize}
    \item il \alert{raggio} è uguale per tutti i cilindri, vengono posizionati
      internamente alla inner rod
      \begin{itemize}
      \item pensare se fornire un \alert{offset} per il raggio da file di
        configurazione
      \end{itemize}
    \end{itemize}
  \item \alert{Calcolo} dei materiali sui cilindri creati
  \end{itemize}
\end{frame}

\end{document}
